\chapter{第二章：AI Agent 行业全景}

AI Agent 的生态在 2026 年已经不再是``几个框架 + 一个模型''的简单组合，而是一套分层技术栈。本章给你一张清晰的行业地图：\textbf{11 类核心工具、主要玩家与企业采用现状。}


\bigskip\hrule\bigskip

\section{2.1 先看总图：Agent 生态是分层的}

站在工程角度，一个完整 Agent 系统通常由六层构成：


\begin{codeblock}
+-----------------------------------------------------------+
| 应用层：客服 / 编码 / 金融分析 / 医疗 / 零售 / 法律 / 教育 |
+-----------------------------------------------------------+
| 执行层：Agent Runtime / Browser Agent / IDE Agent / Memory |
+-----------------------------------------------------------+
| 编排层：Workflow / Scheduling / Tool Routing / A2A / MCP  |
+-----------------------------------------------------------+
| 数据层：RAG / Vector DB / Cache / Session State           |
+-----------------------------------------------------------+
| 运维层：Tracing / Evaluation / Guardrails / Policy        |
+-----------------------------------------------------------+
| 模型层：GPT / Claude / Gemini / Llama / Qwen / ERNIE      |
+-----------------------------------------------------------+
\end{codeblock}

不同公司、不同岗位关注的层次并不一样：模型公司占模型层，框架公司占执行与编排层，企业平台团队关注治理，应用团队关心业务 ROI。


\bigskip\hrule\bigskip

\section{2.2 2026 年 11 类 Agentic AI 工具版图}

根据 StackOne 在 2026 年对 \textbf{120+ Agentic AI 工具}的映射，主流生态大致可分成 11 类。


\subsection{1）基础模型（Foundation Models）}

代表选手包括 GPT-4o / GPT-4.1、Claude 4.x、Gemini 2.x / 3.x、Llama、Qwen、文心（ERNIE）。这一层决定推理、工具选择、长上下文和多模态能力。工程选型时重点看工具调用稳定性、结构化输出、长上下文成本、延迟和合规部署。


\subsection{2）Agent 框架（Agent Frameworks）}

代表选手包括 LangChain / LangGraph、CrewAI、AutoGen、Dify、Coze。它们解决的是``Agent 怎么组织起来''，包括状态管理、工具接口、多步流程、多 Agent 协作和人工审批。


\subsection{3）编排平台（Orchestration Platforms）}

这一层比框架更偏生产环境，关注任务调度、重试超时、审批、工作流版本化和成本控制。很多企业会叠加自己的编排平台。


\subsection{4）向量数据库（Vector Databases）}

代表选手包括 Pinecone、Weaviate、Milvus、FAISS、Chroma。它们是 RAG 的底层基础设施，但别把向量数据库误解成``自动记忆''。真正决定效果的还有 chunking、embedding、元数据过滤、混合检索和 reranking。


\begin{codeblock}
def retrieve(query: str, top_k: int = 3):
    return ["doc-1", "doc-2", "doc-3"][:top_k]


context = retrieve("退款多久到账")
print(context)
\end{codeblock}

\subsection{5）可观测性与评估（Observability \& Evaluation）}

代表选手是 LangSmith、Weights \& Biases、Arize。重点不是``让日志更好看''，而是回答：哪一步失败、哪类任务成功率最低、新版本是否退化、成本和延迟是否可控。


\subsection{6）Agent 协议（Agent Protocols）}

最关键的是 MCP（Model Context Protocol）和 A2A（Agent-to-Agent）。MCP 解决`\texttt{模型或 Agent 如何接工具和资源''，A2A 解决}`不同 Agent 如何协作和传递任务''。协议成熟意味着生态从点对点适配走向标准化接入。


\subsection{7）记忆系统（Memory Systems）}

记忆通常分三层：短期记忆、长期记忆、工作记忆。真正的难点是``什么时候存、什么时候取、什么时候忘''，而不是简单存文本。


\subsection{8）代码/IDE Agent（Code/IDE Agents）}

代表产品包括 Claude Code、Cursor、GitHub Copilot、Windsurf。它们的核心能力是读仓库、搜代码、改文件、跑测试。这是转行者最容易做作品集的一类方向。


\subsection{9）浏览器 Agent（Browser Agents）}

这类 Agent 通过网页环境感知和操作 UI，适合表单填写、运营后台自动化、网页调研和数据抓取。难点在于网页变化频繁，因此依赖重试、回退和审批。


\subsection{10）企业集成平台（Enterprise Integration Platforms）}

企业是否愿意采购 Agent，很大程度上取决于它能不能接进现有系统，如 CRM、ERP、工单系统、邮件系统、数据仓库和 IAM。没有连接器（connectors），Agent 再聪明也只能``会说不会做''。


\subsection{11）安全与护栏（Security \& Guardrails）}

随着 Agent 获得更多操作权限，安全层的重要性快速上升，重点包括 prompt injection 防护、越权控制、敏感操作审批、secret 扫描、审计日志和输出过滤。


一句话总结：\textbf{模型提供智能，框架负责执行，平台负责治理，数据负责知识，观测与安全保证可上线。}


\bigskip\hrule\bigskip

\section{2.3 一张更完整的 ASCII 生态地图}

\begin{codeblock}
                              +----------------------+
                              |  Industry Use Cases  |
                              | CS / Coding / Fin /  |
                              | Healthcare / Legal   |
                              +----------+-----------+
                                         |
        +--------------------+-----------+-----------+--------------------+
        |                    |                       |                    |
        v                    v                       v                    v
 +--------------+   +---------------+     +----------------+   +----------------+
 | IDE Agents   |   | BrowserAgent  |     | Enterprise App |   | Memory Systems |
 | Copilot etc. |   | UI Automation |     | Support/Sales  |   | Session/Long   |
 +------+-------+   +-------+-------+     +--------+-------+   +--------+-------+
        \                   |                      /                    /
         \                  |                     /                    /
          v                 v                    v                    v
      +--------------------------------------------------------------------+
      | Agent Frameworks / Runtime / Orchestration                         |
      | LangChain | CrewAI | AutoGen | Dify | Coze | Workflow | Scheduling |
      +------------------------------+-------------------------------------+
                                     |
                                     v
      +--------------------------------------------------------------------+
      | Protocols & Enterprise Integration                                 |
      | MCP | A2A | Tool Registry | Connectors | IAM | Approval            |
      +------------------------------+-------------------------------------+
                                     |
                                     v
      +--------------------------------------------------------------------+
      | Retrieval / Vector / Knowledge                                     |
      | Pinecone | Weaviate | Milvus | FAISS | Chroma | Cache | RAG        |
      +------------------------------+-------------------------------------+
                                     |
                                     v
      +--------------------------------------------------------------------+
      | Observability / Evaluation / Security                              |
      | LangSmith | W&B | Arize | Tracing | Guardrails | Red Teaming       |
      +------------------------------+-------------------------------------+
                                     |
                                     v
      +--------------------------------------------------------------------+
      | Foundation Models                                                  |
      | GPT | Claude | Gemini | Llama | Qwen | ERNIE                       |
      +--------------------------------------------------------------------+
\end{codeblock}

\bigskip\hrule\bigskip

\section{2.4 主要玩家深度观察}

\subsection{OpenAI}

OpenAI 的强项在于模型能力、API 完整度和开发者心智。


\subsection{Anthropic}

Anthropic 在 2025-2026 年最强的心智是安全、推理质量、coding 场景。Claude 系列在代码理解、仓库导航、工具调用稳定性方面很突出。


\subsection{Google}

Google 的优势是``模型 + 云 + 多模态''。


\subsection{Meta}

Meta 的核心影响力来自 Llama 开源生态。它带来私有化部署、低成本试验和二次微调空间，对强调数据主权的企业有吸引力。


\subsection{Microsoft}

Microsoft 的优势是企业入口：GitHub、Office、Azure、Copilot 家族，再加上成熟的安全与治理能力。很多企业内部 Agent 项目都会优先考虑微软生态。


\subsection{中国玩家：百度、阿里、字节跳动}

中国市场的重要特点是本地模型、云和企业协作工具耦合更紧：\textbf{百度}偏搜索与企业服务，\textbf{阿里巴巴}依托 Qwen 与阿里云，\textbf{字节跳动}产品化速度快。对中文场景求职者来说，本地部署、中文效果、飞书/钉钉/企业微信集成往往更重要。


\bigskip\hrule\bigskip

\section{2.5 企业采用现状：哪些行业已经跑起来了}

\subsection{金融}

典型场景包括财报与研究摘要、合规检索、客服辅助、风险初筛。金融行业愿意投入，因为文档密集、人工分析昂贵。


\subsection{医疗}

典型场景包括分诊、病历摘要、患者沟通辅助、保险理赔文档整理。医疗行业上 Agent 的前提不是`\texttt{更聪明''，而是}`可追踪、可审核、人类最终责任人''。


\subsection{零售与电商}

典型场景包括智能客服、商品描述生成、商家运营助手、库存与补货建议、营销分析。零售之所以是高 ROI 行业，是因为订单量大、重复任务多。


\subsection{法律}

典型场景包括合同条款提取、版本对比、案例检索、初稿生成。法律 Agent 最适合压缩前置机械劳动，而不是替代律师判断。


\subsection{教育}

典型场景包括作业讲解、个性化辅导、备课辅助、学习路径推荐。教育适合长期记忆与多模态能力落地。


\bigskip\hrule\bigskip

\section{2.6 作为转行者，如何看待这个生态}

不要试图把所有工具都学一遍。更有效的方法是问自己五个问题：


\begin{enumerate}[leftmargin=2em]
\item \textbf{它解决的是哪一层的问题？} 模型、检索、框架、编排、评估还是安全？
\item \textbf{它有没有明确业务闭环？} 是能提高工单处理率，还是只会做演示？
\item \textbf{它是否需要工程深度？} 权限、重试、日志、回归测试越多，岗位壁垒通常越强。
\item \textbf{它能否与你原有技能叠加？} 后端适合工具与平台，前端适合 Agent UI，DevOps 适合推理与部署，QA 适合评估与红队测试。
\item \textbf{它能不能做成作品集？} 客服 Agent、编码 Agent、评估平台，都比``会聊天的小机器人''更有说服力。
\end{enumerate}

读完这一章，你至少要形成一个结构：\textbf{Foundation Model 并不等于 Agent；价值主要在模型之上的执行、编排、知识、评估和治理。}


\section{本章要点}

\begin{itemize}[leftmargin=2em]
\item 2026 年 Agentic AI 已形成\textbf{分层明确的生态系统}，不是单一模型或单一框架能解释的。
\item StackOne 的 2026 版图显示，市场上已有 \textbf{11 类、120+ 工具}，说明行业进入快速分工期。
\item 11 类核心类别包括：\textbf{基础模型、Agent 框架、编排平台、向量数据库、可观测性与评估、协议、记忆系统、代码/IDE Agent、浏览器 Agent、企业集成平台、安全与护栏}。
\item OpenAI、Anthropic、Google、Meta、Microsoft 与中国玩家（百度、阿里、字节等）分别在模型、平台、企业入口、开源和本地化部署上形成不同优势。
\item 金融、医疗、零售、法律、教育都已出现较清晰的 Agent 落地场景。
\item 对转行者最重要的，不是记住所有产品名，而是理解每一层解决什么问题，并押注\textbf{有业务闭环、工程深度高}的方向。
\end{itemize}

\section{延伸阅读}

\begin{itemize}[leftmargin=2em]
\item StackOne：2026 Agentic AI Tools Landscape
\item LangChain / LangGraph、CrewAI、AutoGen、Dify、Coze 官方文档
\item Pinecone、Milvus、Weaviate、Chroma 官方文档
\item LangSmith、Weights \& Biases、Arize 关于 LLM 评估的资料
\item MCP 与 A2A 协议规范
\item OpenAI、Anthropic、Google Cloud、Microsoft Azure 的企业 Agent 文档
\end{itemize}
